% LaTeX Document: Strategic Analysis of 2018 AMC 12A Problem 10
% IMPORTANT: Use LaTeX commands for all mathematical symbols, NOT unicode characters

\documentclass[]{article}
\usepackage[utf8]{inputenc}
\usepackage{amsmath, amssymb, amsthm}
\usepackage{geometry}
\usepackage{xcolor}
\usepackage[most]{tcolorbox}
\usepackage{enumitem}
\usepackage{fancyhdr}

% Page setup
\geometry{margin=1in}
\pagestyle{fancy}
\fancyhf{}
\rhead{AMC12/COMC Problem Analysis}
\lhead{\today}

% Color definitions
\definecolor{problemconfig}{RGB}{230, 242, 255}
\definecolor{strategic}{RGB}{255, 230, 242}
\definecolor{tactical}{RGB}{242, 255, 230}
\definecolor{techniques}{RGB}{255, 242, 230}
\definecolor{verification}{RGB}{255, 230, 230}
\definecolor{solution}{RGB}{230, 255, 255}
\definecolor{competition}{RGB}{242, 242, 255}
\definecolor{gold}{RGB}{218, 165, 32}

% Box styles
\newtcolorbox{problemconfigbox}{
    colback=problemconfig,
    colframe=blue,
    title=Problem Configuration Analysis,
    fonttitle=\bfseries,
    arc=3pt,
    boxrule=1pt,
    breakable
}

\newtcolorbox{strategicbox}{
    colback=strategic,
    colframe=purple,
    title=Strategic Framework Application,
    fonttitle=\bfseries,
    arc=3pt,
    boxrule=1pt,
    breakable
}

\newtcolorbox{tacticalbox}{
    colback=tactical,
    colframe=green,
    title=Tactical Methodology Selection,
    fonttitle=\bfseries,
    arc=3pt,
    boxrule=1pt,
    breakable
}

\newtcolorbox{techniquesbox}{
    colback=techniques,
    colframe=orange,
    title=Conversion Techniques Application,
    fonttitle=\bfseries,
    arc=3pt,
    boxrule=1pt,
    breakable
}

\newtcolorbox{verificationbox}{
    colback=verification,
    colframe=red,
    title=Verification Strategy Design,
    fonttitle=\bfseries,
    arc=3pt,
    boxrule=1pt,
    breakable
}

\newtcolorbox{solutionbox}{
    colback=solution,
    colframe=cyan,
    title=Solution Roadmap,
    fonttitle=\bfseries,
    arc=3pt,
    boxrule=1pt,
    breakable
}

\newtcolorbox{competitionbox}{
    colback=competition,
    colframe=purple,
    title=Competition Strategy Integration,
    fonttitle=\bfseries,
    arc=3pt,
    boxrule=1pt,
    breakable
}

\newtcolorbox{keyinsightbox}{
    colback=yellow!20,
    colframe=gold,
    title=Key Insight,
    fonttitle=\bfseries,
    arc=3pt,
    boxrule=2pt,
    leftrule=5pt,
    breakable
}

\newtcolorbox{warningbox}{
    colback=red!10,
    colframe=red!50,
    title=Warning: Common Pitfall,
    fonttitle=\bfseries,
    arc=3pt,
    boxrule=2pt,
    leftrule=5pt,
    breakable
}

\newtcolorbox{tipbox}{
    colback=green!10,
    colframe=green!50,
    title=Pro Tip,
    fonttitle=\bfseries,
    arc=3pt,
    boxrule=2pt,
    leftrule=5pt,
    breakable
}

\begin{document}

\title{AMC12/COMC Strategic Problem Analysis\\2018 AMC 12A Problem 10}
\author{Competition Math Framework}
\date{\today}
\maketitle

\section*{Problem Statement}
How many ordered pairs $(x,y)$ of real numbers satisfy the following system of equations?
\begin{align*}
x + 3y &= 3 \\
\left| |x| - |y| \right| &= 1
\end{align*}

$\textbf{(A) } 1 \qquad \textbf{(B) } 2 \qquad \textbf{(C) } 3 \qquad \textbf{(D) } 4 \qquad \textbf{(E) } 8$

\begin{problemconfigbox}
\subsection*{A. Problem Classification}
\begin{itemize}[leftmargin=*]
    \item \textbf{Problem Type}: To Find (counting ordered pairs)
    \item \textbf{Difficulty Band}: Early-to-Mid-band (Problem 10 on AMC 12A)
    \item \textbf{Competition Context}: AMC12 (multiple choice, 6 points)
    \item \textbf{Mathematical Domain}: Algebra (systems of equations with absolute values)
    \item \textbf{Estimated Time}: 3-5 minutes for mid-band problem
\end{itemize}

\subsection*{B. Problem Structure Identification}
\begin{itemize}[leftmargin=*]
    \item \textbf{Given Information}: 
    \begin{itemize}
        \item Linear constraint: $x + 3y = 3$
        \item Nested absolute value constraint: $||x| - |y|| = 1$
    \end{itemize}
    \item \textbf{Constraints}: 
    \begin{itemize}
        \item Real numbers (continuous domain)
        \item Must satisfy both equations simultaneously
        \item Absolute value structure creates case-based analysis
    \end{itemize}
    \item \textbf{Objective}: Count the number of ordered pairs $(x,y)$ satisfying both equations
    \item \textbf{Hidden Assumptions}: 
    \begin{itemize}
        \item Nested absolute values require careful case analysis
        \item Sign considerations for $x$ and $y$
        \item Some cases may yield duplicate solutions
    \end{itemize}
    \item \textbf{Answer Choice Leverage}: 
    \begin{itemize}
        \item Small integer answers (1-8) suggest complete enumeration is feasible
        \item Answer (E) 8 suggests $2^3$ structure (3 binary choices)
        \item Answer (C) 3 is "middle ground" - often correct in AMC12
    \end{itemize}
\end{itemize}

\subsection*{C. Strategic Orientation}
\begin{itemize}[leftmargin=*]
    \item \textbf{Hypothesis}: Given $x + 3y = 3$ and $||x| - |y|| = 1$
    \item \textbf{Conclusion}: Find the count of ordered pairs $(x,y)$
    \item \textbf{Notation}: Use $(x,y)$ for ordered pairs; consider signs separately
    \item \textbf{Intuitive Approach}: 
    \begin{itemize}
        \item Use first equation to express $x = 3 - 3y$
        \item Substitute into second equation
        \item Break into cases based on absolute value signs
    \end{itemize}
    \item \textbf{Cross-Domain Potential}: 
    \begin{itemize}
        \item Geometric interpretation: intersection of line with absolute value curves
        \item Visualization could provide insight into solution count
    \end{itemize}
\end{itemize}
\end{problemconfigbox}

\begin{strategicbox}
\subsection*{A. Psychological Strategies}
\begin{itemize}[leftmargin=*]
    \item \textbf{Mental Toughness}: Absolute value problems require patience with casework - commit to systematic approach rather than getting overwhelmed
    \item \textbf{Creativity Cultivation}: Consider both algebraic substitution and geometric visualization
    \item \textbf{Iterative Refinement}: After finding solutions, verify they truly satisfy both original equations (not just simplified forms)
\end{itemize}

\subsection*{B. Investigation Strategies}
\begin{itemize}[leftmargin=*]
    \item \textbf{Startup Orientation}: Recognize this as an absolute value casework problem - prepare to handle multiple cases systematically
    \item \textbf{Get Your Hands Dirty}: Substitute $x = 3-3y$ immediately to reduce to one variable
    \item \textbf{Penultimate Step}: After finding all $(x,y)$ values, the final step is counting distinct ordered pairs
    \item \textbf{Wishful Thinking}: "What if the nested absolute value simplifies nicely?" - This guides us to consider when $|x|-|y|$ is positive vs negative
    \item \textbf{Make It Easier}: Work with the substituted single-variable equation first
    \item \textbf{Conjecture via Small Cases}: After substitution, we get $||3-3y| - |y|| = 1$ - test specific values like $y=0, 1, -1$ to understand behavior
    \item \textbf{Visualization}: The line $x+3y=3$ intersects curves defined by $||x|-|y||=1$ - imagining this helps predict solution count
\end{itemize}

\subsection*{C. Argument Construction Strategy}
\begin{itemize}[leftmargin=*]
    \item \textbf{Direct Proof}: Use algebraic case analysis to enumerate all solutions
    \item \textbf{Exhaustive Casework}: Break problem into cases:
    \begin{enumerate}
        \item Outer absolute value: $|x| - |y| = 1$ or $|x| - |y| = -1$
        \item For each, consider signs: $x \geq 0$ or $x < 0$, and $y \geq 0$ or $y < 0$
    \end{enumerate}
\end{itemize}

\subsection*{D. Cross-Domain Translation}
\begin{itemize}[leftmargin=*]
    \item \textbf{Algebraic Approach}: Substitute and solve cases systematically
    \item \textbf{Geometric Visualization}: Plot the line and absolute value curves to see intersection points
    \item \textbf{Primary Method}: Algebraic (more reliable for exact counting)
    \item \textbf{Secondary Check}: Geometric intuition for sanity check
\end{itemize}
\end{strategicbox}

\begin{tacticalbox}
\subsection*{A. Core Mathematical Tactics}
\begin{itemize}[leftmargin=*]
    \item \textbf{Primary Tactic}: \textit{Substitution} - Reduce two-variable system to single variable
    \item \textbf{Secondary Tactic}: \textit{Case Analysis} - Handle absolute value conditions systematically
    \item \textbf{Supporting Tactics}:
    \begin{itemize}
        \item Algebraic manipulation (solving linear/quadratic equations)
        \item Sign analysis (determining when expressions are positive/negative)
        \item Solution verification (checking each candidate against original equations)
    \end{itemize}
\end{itemize}

\subsection*{B. Domain-Specific Techniques}
\begin{itemize}[leftmargin=*]
    \item \textbf{Absolute Value Handling}:
    \begin{itemize}
        \item $|A| = B$ means $A = B$ or $A = -B$
        \item Nested absolute values require careful ordering: work from inside out
        \item Track sign conditions throughout
    \end{itemize}
    \item \textbf{System Solving}:
    \begin{itemize}
        \item Linear equation allows direct substitution
        \item Convert to single-variable problem
    \end{itemize}
\end{itemize}

\subsection*{C. Crossover Techniques}
\begin{itemize}[leftmargin=*]
    \item \textbf{Algebraic $\leftrightarrow$ Geometric}: Line-curve intersection interpretation
    \item \textbf{Numeric $\leftrightarrow$ Symbolic}: Test specific values to guide case analysis
\end{itemize}
\end{tacticalbox}

\begin{techniquesbox}
\subsection*{A. Recognition Index}
\textbf{Signature Pattern Identified}: System with linear equation + nested absolute value equation

\textbf{High-Probability Techniques}:
\begin{itemize}[leftmargin=*]
    \item \textbf{Substitution Method} (90\% applicability): Linear equation immediately suggests expressing one variable in terms of the other
    \item \textbf{Casework Analysis} (95\% applicability): Nested absolute values mandate case-by-case breakdown
\end{itemize}

\subsection*{B. Technique Selection}
\textbf{Selected Approach}: Substitution followed by nested absolute value casework

\textbf{Justification}:
\begin{itemize}[leftmargin=*]
    \item Linear equation is "free" - trivial to solve for $x$ in terms of $y$
    \item Reduces two-variable problem to one-variable problem
    \item Absolute value structure requires cases anyway
    \item Alternative (solving for $y$ first) would work but is messier due to coefficient 3
\end{itemize}

\subsection*{C. Integration Strategy}
\textbf{Step-by-step synthesis}:
\begin{enumerate}
    \item Substitute $x = 3-3y$ from first equation into second equation
    \item Obtain $||3-3y| - |y|| = 1$
    \item Apply outer absolute value: two cases
    \begin{itemize}
        \item Case A: $|3-3y| - |y| = 1$
        \item Case B: $|3-3y| - |y| = -1$
    \end{itemize}
    \item For each case, apply inner absolute values based on sign of $3-3y$ and $y$
    \item Solve resulting linear equations for $y$
    \item Compute corresponding $x$ values
    \item Verify and count distinct solutions
\end{enumerate}

\subsection*{D. Potential Conflicts and Resolutions}
\begin{itemize}[leftmargin=*]
    \item \textbf{Conflict}: Multiple case paths might yield same solution
    \item \textbf{Resolution}: Track all solutions, then count distinct ordered pairs at end
    \item \textbf{Conflict}: Sign conditions might be inconsistent with derived values
    \item \textbf{Resolution}: Verify each solution satisfies its case conditions before accepting
\end{itemize}
\end{techniquesbox}

\begin{verificationbox}
\subsection*{A. Verification Level Selection}
\textbf{Selected Level}: Level 4 (Moderate-High)

\textbf{Decision Tree Justification}:
\begin{itemize}[leftmargin=*]
    \item \textbf{Problem Difficulty}: Mid-band (P10) - not trivial, deserves careful checking
    \item \textbf{Computation Complexity}: Multiple cases with potential for sign errors
    \item \textbf{Answer Impact}: Counting problem - off by one error is common
    \item \textbf{Time Available}: 3-5 minute problem allows 30-60 seconds for verification
\end{itemize}

\subsection*{B. Verification Methods}
\begin{enumerate}
    \item \textbf{Direct Substitution}: Plug each $(x,y)$ candidate into both original equations
    \item \textbf{Case Consistency Check}: Verify sign assumptions hold for each solution
    \item \textbf{Graphical Sanity Check}: Mentally visualize if solution count makes geometric sense
    \item \textbf{Answer Choice Validation}: Does final count match one of the given options?
\end{enumerate}

\subsection*{C. Confidence Calibration}
\begin{itemize}[leftmargin=*]
    \item \textbf{Target Confidence}: 90-95\% (appropriate for mid-band problem)
    \item \textbf{Confidence Indicators}:
    \begin{itemize}
        \item All solutions satisfy both original equations upon substitution
        \item No duplicate solutions found
        \item Sign conditions consistent with derived values
        \item Final count matches an answer choice
    \end{itemize}
    \item \textbf{Red Flags}:
    \begin{itemize}
        \item Solutions that violate their case assumptions
        \item Suspiciously high solution count (suggests overcounting)
        \item Arithmetic errors in solving linear equations
    \end{itemize}
\end{itemize}
\end{verificationbox}

\begin{solutionbox}
\subsection*{A. Initial Setup (30 seconds)}
From $x + 3y = 3$, express:
\[x = 3 - 3y\]

Substitute into second equation:
\[\left| |3-3y| - |y| \right| = 1\]

\subsection*{B. Case Analysis (3 minutes)}

\textbf{Outer Absolute Value}: Two primary cases

\textbf{Case A}: $|3-3y| - |y| = 1$

\quad \textbf{Subcase A1}: $3-3y \geq 0$ (i.e., $y \leq 1$) and $y \geq 0$
\begin{align*}
(3-3y) - y &= 1\\
3 - 4y &= 1\\
y &= \frac{1}{2}
\end{align*}
Check: $y = \frac{1}{2} \in [0,1]$ \checkmark

Then $x = 3 - 3(\frac{1}{2}) = \frac{3}{2}$

\textbf{Solution 1}: $\left(\frac{3}{2}, \frac{1}{2}\right)$

\quad \textbf{Subcase A2}: $3-3y \geq 0$ and $y < 0$
\begin{align*}
(3-3y) - (-y) &= 1\\
3 - 3y + y &= 1\\
3 - 2y &= 1\\
y &= 1
\end{align*}
Check: $y = 1 \not< 0$ \texttimes \quad (Contradiction - discard)

\quad \textbf{Subcase A3}: $3-3y < 0$ (i.e., $y > 1$) and $y \geq 0$
\begin{align*}
-(3-3y) - y &= 1\\
-3 + 3y - y &= 1\\
2y &= 4\\
y &= 2
\end{align*}
Check: $y = 2 > 1$ \checkmark

Then $x = 3 - 3(2) = -3$

\textbf{Solution 2}: $(-3, 2)$

\quad \textbf{Subcase A4}: $3-3y < 0$ and $y < 0$
\begin{align*}
-(3-3y) - (-y) &= 1\\
-3 + 3y + y &= 1\\
4y &= 4\\
y &= 1
\end{align*}
Check: $y = 1 \not< 0$ and $y = 1 \not> 1$ \texttimes \quad (Contradiction - discard)

\textbf{Case B}: $|3-3y| - |y| = -1$

\quad \textbf{Subcase B1}: $3-3y \geq 0$ and $y \geq 0$
\begin{align*}
(3-3y) - y &= -1\\
3 - 4y &= -1\\
y &= 1
\end{align*}
Check: $y = 1 \in [0,1]$ \checkmark

Then $x = 3 - 3(1) = 0$

\textbf{Solution 3}: $(0, 1)$

\quad \textbf{Subcase B2}: $3-3y \geq 0$ and $y < 0$
\begin{align*}
(3-3y) + y &= -1\\
3 - 2y &= -1\\
y &= 2
\end{align*}
Check: $y = 2 \not< 0$ \texttimes \quad (Contradiction - discard)

\quad \textbf{Subcase B3}: $3-3y < 0$ and $y \geq 0$
\begin{align*}
-(3-3y) - y &= -1\\
-3 + 3y - y &= -1\\
2y &= 2\\
y &= 1
\end{align*}
Check: $y = 1 \not> 1$ \texttimes \quad (Boundary case from A1 - already counted)

\quad \textbf{Subcase B4}: $3-3y < 0$ and $y < 0$
\begin{align*}
-(3-3y) + y &= -1\\
-3 + 3y + y &= -1\\
4y &= 2\\
y &= \frac{1}{2}
\end{align*}
Check: $y = \frac{1}{2} \not< 0$ and $y = \frac{1}{2} \not> 1$ \texttimes \quad (Contradiction - discard)

\subsection*{C. Solution Summary (30 seconds)}
\textbf{Valid ordered pairs}:
\begin{enumerate}
    \item $\left(\frac{3}{2}, \frac{1}{2}\right)$ from Subcase A1
    \item $(-3, 2)$ from Subcase A3
    \item $(0, 1)$ from Subcase B1
\end{enumerate}

\subsection*{D. Verification (1 minute)}
\textbf{Check each solution in both original equations}:

\textbf{Solution 1}: $\left(\frac{3}{2}, \frac{1}{2}\right)$
\begin{itemize}
    \item $x + 3y = \frac{3}{2} + 3 \cdot \frac{1}{2} = \frac{3}{2} + \frac{3}{2} = 3$ \checkmark
    \item $||x| - |y|| = \left|\frac{3}{2} - \frac{1}{2}\right| = |1| = 1$ \checkmark
\end{itemize}

\textbf{Solution 2}: $(-3, 2)$
\begin{itemize}
    \item $x + 3y = -3 + 3 \cdot 2 = -3 + 6 = 3$ \checkmark
    \item $||x| - |y|| = |3 - 2| = |1| = 1$ \checkmark
\end{itemize}

\textbf{Solution 3}: $(0, 1)$
\begin{itemize}
    \item $x + 3y = 0 + 3 \cdot 1 = 3$ \checkmark
    \item $||x| - |y|| = |0 - 1| = |-1| = 1$ \checkmark
\end{itemize}

All three solutions are valid!

\subsection*{E. Pitfall Prevention}
\begin{itemize}[leftmargin=*]
    \item \textbf{Pitfall 1}: Forgetting to check if solutions satisfy case conditions
    \item \textbf{Mitigation}: Explicitly verified each $y$ value against its case assumptions
    \item \textbf{Pitfall 2}: Double-counting solutions that arise from multiple cases
    \item \textbf{Mitigation}: Listed distinct solutions and verified each only once
    \item \textbf{Pitfall 3}: Sign errors in absolute value expansion
    \item \textbf{Mitigation}: Carefully tracked signs in each subcase
\end{itemize}
\end{solutionbox}

\begin{competitionbox}
\subsection*{A. Time Management}
\begin{itemize}[leftmargin=*]
    \item \textbf{Estimated Time}: 4-5 minutes
    \item \textbf{Time Breakdown}:
    \begin{itemize}
        \item Setup and substitution: 30 seconds
        \item Case analysis: 3 minutes
        \item Solution counting: 30 seconds
        \item Verification: 1 minute
    \end{itemize}
    \item \textbf{Actual Pace}: On target for Problem 10 (mid-band)
\end{itemize}

\subsection*{B. Problem Prioritization}
\begin{itemize}[leftmargin=*]
    \item \textbf{Accessibility}: Medium - requires systematic casework but no advanced techniques
    \item \textbf{When to Attempt}: During first pass of mid-band problems (P9-16)
    \item \textbf{Skip Conditions}: If uncomfortable with absolute values, mark and return later
\end{itemize}

\subsection*{C. Risk Management}
\begin{itemize}[leftmargin=*]
    \item \textbf{Confidence Level}: 95\% - All solutions verified against original equations
    \item \textbf{Flag for Review}: No - high confidence with complete verification
    \item \textbf{Guessing Strategy (if needed)}: Answer (C) is often correct for AMC12; geometric intuition suggests 3-4 intersection points
\end{itemize}

\subsection*{D. Abandonment Criteria}
\begin{itemize}[leftmargin=*]
    \item \textbf{Soft Abandon}: If exceeding 7 minutes total
    \item \textbf{Hard Abandon}: If making repeated algebraic errors or case confusion
    \item \textbf{Partial Credit Strategy}: N/A (AMC12 has no partial credit)
\end{itemize}
\end{competitionbox}

\begin{keyinsightbox}
\textbf{Key Insight}: After substitution, the nested absolute value equation reduces to a manageable number of cases (8 subcases total, but many yield contradictions). The linear equation makes substitution trivial, converting the problem to single-variable absolute value cases. The answer is \textbf{exactly 3 solutions}.
\end{keyinsightbox}

\begin{warningbox}
\textbf{Warning}: The most common pitfall is \textbf{failing to verify that solutions satisfy their case conditions}. For example, Subcase A2 gives $y=1$, but we assumed $y < 0$, so this must be discarded. Always check sign assumptions after solving!
\end{warningbox}

\begin{tipbox}
\textbf{Pro Tip}: When dealing with nested absolute values, work systematically through all cases even if some seem unlikely. Use a clear notation system (like "A1, A2, ..." for subcases) to avoid confusion. This organizational strategy prevents missing solutions or double-counting.
\end{tipbox}

\section*{Final Answer}
\boxed{\textbf{(C) } 3}

\section*{Confidence Assessment}
\begin{itemize}[leftmargin=*]
    \item \textbf{Confidence Level}: 95\% - All three solutions independently verified
    \item \textbf{Verification Applied}: Level 4 (Direct substitution + case consistency)
    \item \textbf{Flag for Review}: No - complete case analysis with no contradictions
    \item \textbf{Time Spent}: Approximately 4-5 minutes (optimal for Problem 10)
\end{itemize}

\end{document}